%=====================================================================
\chapter{Administration: The Dashboard, and What to Watch}\label{ch:dashboard}
%=====================================================================
\locator{5}{Part V --- Running and Extending the Shop}

\begin{outcomes}
By the end of this chapter you will be able to:
\begin{tight}
  \item \textbf{Describe} what each dashboard widget counts, and \textbf{state}
        the filter behind it.
  \item \textbf{Explain} why the dashboard's numbers disagree, and
        \textbf{reconcile} them against the underlying table.
  \item \textbf{Distinguish} a figure that describes the shop from a figure
        that would change a decision.
  \item \textbf{Derive} three measurements this course has produced that the
        dashboard does not show, and \textbf{write} the query for each.
  \item \textbf{Design} a panel for a stated shop, and \textbf{justify} every
        number on it.
\end{tight}
\end{outcomes}

\begin{prereq}
\chref{checkout}: order statuses, and status 0. \chref{catalogue}: reading the
query a screen actually runs. \chref{shipping} and \chref{payments}: the costs
that make some numbers worth watching and others not.
\end{prereq}

\begin{keyterms}
dashboard $\bullet$ widget $\bullet$ report $\bullet$ filter $\bullet$ vanity
metric $\bullet$ actionable metric $\bullet$ reconciliation $\bullet$
abandonment rate $\bullet$ conversion $\bullet$ leading and lagging indicator
$\bullet$ denominator
\end{keyterms}

\section{Eight Widgets, and What Each One Counts}

The administrator's first screen carries eight panels. They are extensions like
any other, each switched on at \adminpath{Extensions, Extensions, Dashboard},
each with a width and a sort order.

\rowcolors{2}{white}{lightbg}
\begin{longtable}{@{}L{2.4cm}L{5.6cm}L{6.4cm}@{}}
\toprule
\rowcolor{primary!15}
\textbf{Widget} & \textbf{Shows} & \textbf{Worth knowing} \\
\midrule
\endhead
\texttt{order} & Orders today, yesterday, in total & Counts every order above
status 0 \\
\texttt{sale} & Money, over the same periods & Same filter, different query \\
\texttt{customer} & Registered customers & Not buyers --- registrations \\
\texttt{online} & Visitors active now & \textbf{Empty unless you switch on
tracking} \\
\texttt{activity} & What customers have been doing & \textbf{Also empty by
default} \\
\texttt{chart} & Orders over time & \textbf{Counts only completed statuses} \\
\texttt{map} & Where orders came from & By billing country \\
\texttt{recent} & The last few orders & The most useful of the eight \\
\bottomrule
\end{longtable}

\begin{alertbox}[Two of the eight are permanently empty, and that is a setting]
\texttt{config\_customer\_online} and \texttt{config\_customer\_activity} are
both \textbf{0} in a fresh install, so the \texttt{online} and
\texttt{activity} panels have nothing to show and never will until somebody
changes them.

\smallskip
\textbf{Before switching them on, recall \chref{basket}.} Both work by writing
a row for each visitor's activity --- so turning them on adds a database write
to page views that previously had none, on a shop whose session table already
accumulated 645 rows without a single customer. The panels are not free; they
are paid for in writes on your busiest pages.

\smallskip
\textbf{Decide what you would do differently if you knew.} If the answer is
``nothing'', leave them off. That question is the whole of this chapter.
\end{alertbox}

\section{The Numbers Disagree, and Nothing Says So}

The four headline figures are produced by different queries. This is not a
fault --- each filter is right for its own purpose --- but the screen presents
them side by side with no indication that they are answering different
questions.

\begin{worked}
\textbf{A known set of orders.} Produced by \texttt{code/m17\_dashboard.py}:
five orders paid for through the checkout, two of which were then moved to
\uiel{Complete}, plus one abandoned at the payment button.

\smallskip
\textbf{What is actually in \texttt{oc\_order}:}

\smallskip
\begin{tabular}{@{}L{3.0cm}L{2.6cm}L{3.4cm}L{4.4cm}@{}}
\textbf{Status} & \textbf{Orders} & \textbf{Value} & \\
0 & 1 & 854.40 & \emph{not a status at all} \\
1 & 3 & 2{,}563.20 & Pending \\
5 & 2 & 1{,}708.80 & Complete \\
\end{tabular}

\smallskip
\textbf{What each part of the dashboard reports of exactly that data:}

\smallskip
\begin{tabular}{@{}L{5.1cm}L{2.4cm}L{3.0cm}L{2.8cm}@{}}
& \textbf{Orders} & \textbf{Value} & \textbf{Filter} \\
Orders widget & 5 & 4{,}272.00 & status $>0$ \\
Sales widget & 5 & 4{,}272.00 & status $>0$ \\
\textbf{The chart} & \textbf{2} & \textbf{1{,}708.80} & status in (5,\,3) \\
The table itself & 6 & 5{,}126.40 & everything \\
\end{tabular}

\smallskip
\textbf{Three different answers to ``how many orders'', all correct.} Six in
the table, five on the widgets, two on the chart --- and a shopkeeper looking
at the screen sees ``5'' above a chart whose bars add up to 2, with nothing
explaining the difference.

\smallskip
\textbf{Both gaps are information, once you know they are there.}

\smallskip
\textbf{Six minus five is the abandonment.} One order in six reached the
payment button and stopped. That is the single most useful number a small shop
can compute, and it exists only as the difference between two figures the
dashboard shows separately.

\smallskip
\textbf{Five minus two is the work not yet done.} Three orders are sitting at
Pending: paid for, not shipped. On a Monday morning that is the day's job. On a
Friday it is a problem.

\smallskip
\textbf{And the chart's filter is a \emph{setting}}, not a constant:
\texttt{config\_complete\_status} is \texttt{["5","3"]} here, so Shipped counts
as complete. Change it and every historical bar changes with it.
\end{worked}

\begin{conceptbox}[How to read any dashboard you did not build]
\begin{tightnum}
  \item \textbf{For each number, ask what it counts and what it excludes.} The
        answer is in a query, and you now know how to find it.
  \item \textbf{Reconcile against the table.} If the parts do not add up to the
        whole, the difference has a name --- and the name is usually the
        interesting number.
  \item \textbf{Find which filters are settings.} Those can change under you
        without any code changing.
  \item \textbf{Check the denominator.} ``Conversion'' means nothing until you
        know whether the bottom of the fraction is visitors, sessions, baskets
        or people.
\end{tightnum}
\end{conceptbox}

\section{What the Dashboard Does Not Show}

Every measurement in this course that would change a decision is absent from
the administrator's first screen. That is not an accident: the dashboard shows
what is cheap to count, and those are not the same things.

\begin{examplebox}[Three panels this course has already earned]
\begin{tightnum}
  \item \textbf{Abandonment rate.} From \chref{checkout}: orders at status 0
        against orders above it, by day.
        \smallskip
        \texttt{SELECT DATE(date\_added) d,}\\
        \texttt{\ \ SUM(order\_status\_id=0) abandoned,}\\
        \texttt{\ \ SUM(order\_status\_id>0) placed}\\
        \texttt{FROM oc\_order GROUP BY d ORDER BY d DESC;}
        \smallskip
        A shop with a rising ratio has a checkout problem, a payment problem or
        a price problem, and knows to look before customers tell it.
  \item \textbf{Products that cannot be delivered profitably.} From
        \chref{shipping}: enabled products with a weight of zero, which under
        weight-based shipping are quoted free delivery.
        \smallskip
        \texttt{SELECT COUNT(*) FROM oc\_product}\\
        \texttt{WHERE status=1 AND weight=0;}
        \smallskip
        A number that should be zero and silently is not.
  \item \textbf{Orders standing still.} From this chapter: how long orders have
        been at Pending.
        \smallskip
        \texttt{SELECT order\_status\_id,}\\
        \texttt{\ \ COUNT(*) n,}\\
        \texttt{\ \ MIN(date\_added) oldest}\\
        \texttt{FROM oc\_order WHERE order\_status\_id>0}\\
        \texttt{GROUP BY order\_status\_id;}
        \smallskip
        The widgets say how many orders arrived. Nothing says how many are
        waiting, or for how long --- which is the only one of the three a
        shopkeeper can act on this morning.
\end{tightnum}

\smallskip
All three are one query each, and all three are \chref{modules}'s subject: a
dashboard widget is the smallest useful module you can write.
\end{examplebox}

\dsfig{%
\begin{tikzpicture}[
  box/.style={rounded corners=3pt, draw=#1, line width=1.0pt, fill=#1!8,
              text=#1!60!black, font=\bfseries\scriptsize, align=center,
              text width=3.3cm, minimum height=1.4cm},
  note/.style={font=\scriptsize, text=neutral, align=center, text width=3.6cm}
]
  \node[box=neutral]   (a) at (0,0)    {Describes the shop\\[2pt]
       {\scriptsize\mdseries orders today, total customers, visitors now}};
  \node[box=secondary] (b) at (4.3,0)  {Tracks a trend\\[2pt]
       {\scriptsize\mdseries orders this week against last}};
  \node[box=greenhl]   (c) at (8.6,0)  {Changes a decision\\[2pt]
       {\scriptsize\mdseries orders waiting since Friday}};

  \node[note] at (0,-1.5)   {Nice to see once. Nobody acts on it twice.};
  \node[note] at (4.3,-1.5) {Worth watching. Acted on next month.};
  \node[note, text=greenhl] at (8.6,-1.5)
       {\textbf{Acted on today.} Put these at the top.};

  \draw[-{Stealth[length=2.2mm]}, primary!50, line width=0.9pt]
       (-1.9,-2.5) -- (10.5,-2.5);
  \node[font=\scriptsize\itshape, text=primary] at (4.3,-2.9)
       {increasing usefulness --- and decreasing likelihood of being on the
        default dashboard};
\end{tikzpicture}}%
{A test for every number you are thinking of putting on a screen. All eight
shipped widgets sit at the left-hand end.}{metricsladder}

\begin{pitfall}
\begin{tight}
  \item \textbf{Comparing two dashboard numbers that have different filters.}
        Five orders and a chart showing two is not a bug report.
  \item \textbf{Reporting revenue from a table count} without excluding status
        0.
  \item \textbf{Switching on the online and activity panels} without noticing
        the write they add to every page view.
  \item \textbf{Treating \texttt{config\_complete\_status} as a constant.} It
        is a setting, and changing it rewrites history.
  \item \textbf{Adding a number nobody would act on.} Every panel costs a query
        on the busiest screen in the admin panel.
  \item \textbf{Quoting a conversion rate without naming its denominator.}
  \item \textbf{Watching totals rather than ages.} ``How many orders'' is a
        report; ``how long has this one been waiting'' is a job.
\end{tight}
\end{pitfall}

\begin{lab}[Make the dashboard disagree with itself, then fix the gap]
\begin{lstlisting}[language=bash]
# --- 1. Build a known set of orders ---------------------------------
cd code
python m17_dashboard.py

# --- 2. Now look at the admin dashboard -----------------------------
# Admin: Dashboard. Write down the four headline numbers and what the
# chart shows. Compare with what step 1 printed.

# --- 3. Find the filter behind each number --------------------------
docker compose exec shop sh -c \
  'grep -n "order_status_id" admin/model/sale/order.php | head -3'
# Then do the same for the chart's model. Three different filters.

# --- 4. Account for both gaps ---------------------------------------
# 6 in the table, 5 on the widgets, 2 on the chart. Name each gap.

# --- 5. Change a setting and watch history move ---------------------
# Admin: System > Settings > Edit > Option. Add "Processing" to the
# complete order statuses. Reload the dashboard. Which numbers moved?

# --- 6. Write the three queries the dashboard does not have ---------
# Abandonment by day; enabled products with zero weight; orders by
# status with the age of the oldest. Run each one.

# --- 7. Decide what you would put on the screen ---------------------
# Six panels, no more. For each: the number, the query, and one
# sentence saying what you would DO if it moved.

# --- 8. Clean up ----------------------------------------------------
python m17_dashboard.py --clean
\end{lstlisting}

\smallskip
\textbf{What to notice.} Step 7 is harder than it looks. Most proposed panels
fail the ``what would you do'' test, and the ones that survive are rarely the
ones that were easiest to count.
\end{lab}

\begin{checkpoint}
\begin{tightnum}
  \item The table holds 6 orders worth 5{,}126.40; the widgets report 5 orders
        worth 4{,}272.00; the chart shows 2. Explain each gap, say which of the
        three figures you would put in front of the shop's owner, and what you
        would label it.
  \item A shop owner says ``our conversion rate is 3\%''. List the questions you
        would ask before accepting that number, and explain why two shops both
        reporting 3\% may not be comparable at all.
  \item Propose a dashboard panel that would have caught \chref{checkout}'s
        broken guest checkout --- where every page returned 200 and no order
        could be placed. State the query, how often it runs, and the condition
        that raises the alarm.
\end{tightnum}
\end{checkpoint}

\begin{chaptersummary}
\begin{tight}
  \item Eight widgets ship. Two of them --- online and activity --- are empty
        until tracking is switched on, and that tracking costs a database write
        on page views that previously had none.
  \item Worked: the same six orders were reported as 6, 5 and 2 by the table,
        the widgets and the chart. All three are correct; none is labelled with
        its filter.
  \item The gaps are the information. Table minus widgets is the abandonment
        rate; widgets minus chart is the work not yet done.
  \item The chart's filter is \texttt{config\_complete\_status}, a setting ---
        so changing it rewrites every historical bar.
  \item To read any dashboard: ask what each number excludes, reconcile against
        the table, find which filters are settings, and check the denominator.
  \item A number earns its place on a screen only if somebody would do
        something differently on seeing it move. Most shipped widgets describe;
        few change a decision.
  \item Three panels this course has already earned, one query each:
        abandonment by day, enabled products with zero weight, and orders by
        status with the age of the oldest.
  \item Totals are a report; ages are a job. Prefer the number that names
        today's work.
\end{tight}
\end{chaptersummary}

\begin{reviewq}
  \item \textbf{(MCQ)} The dashboard chart counts orders whose status is:
        (a)~greater than zero; (b)~in \texttt{config\_complete\_status};
        (c)~any status; (d)~Pending only.
  \item \textbf{(MCQ)} The difference between the order count in the table and
        the order count on the widgets is: (a)~a bug; (b)~refunded orders;
        (c)~orders abandoned at the payment button; (d)~orders from other
        stores.
  \item \textbf{(Short)} Explain why switching on the ``customers online''
        panel is not free, referring to a finding from \chref{basket}.
  \item \textbf{(Short)} Give the test that decides whether a number belongs on
        a dashboard, and apply it to ``total registered customers''.
  \item \textbf{(Short)} Why does changing \texttt{config\_complete\_status}
        alter charts of months that have already happened?
  \item \textbf{(Applied)} A shop's orders widget has been flat for three weeks
        while its sales widget has fallen 20\%. Give three distinct
        explanations consistent with both, say what you would query to
        distinguish them, and state which explanation you would investigate
        first and why.
  \item \textbf{(Applied)} Design the dashboard for a shop taking 40 orders a
        day on cash on delivery. Give six panels, the query behind each, how
        often it refreshes, and for every one a sentence beginning ``if this
        moves, we will\dots''. Justify anything you left off that the shipped
        dashboard includes.
\end{reviewq}
